%!TEX TS-program = lualatex
%!TEX encoding = UTF-8 Unicode

\documentclass[t]{beamer}

%%%% HANDOUTS For online Uncomment the following four lines for handout
%\documentclass[t,handout]{beamer}  %Use this for handouts.
%\usepackage{handoutWithNotes}
%\pgfpagesuselayout{3 on 1 with notes}[letterpaper,border shrink=5mm]
%	\setbeamercolor{background canvas}{bg=black!5}

\includeonlylecture{student}

%%% Including only some slides for students.
%%% Uncomment the following line. For the slides,
%%% use the labels shown below the command.

%% For students, use \lecture{student}{student}
%% For mine, use \lecture{instructor}{instructor}

% FONTS
\usepackage{fontspec}
\def\mainfont{Linux Biolinum O}
\setmainfont[Ligatures={Common,TeX}, Contextuals={NoAlternate}, BoldFont={* Bold}, ItalicFont={* Italic}, Numbers={OldStyle}]{\mainfont}
\setsansfont[Ligatures={Common,TeX}, Scale=MatchLowercase, Numbers=OldStyle, BoldFont={* Bold}, ItalicFont={* Italic},]{Linux Biolinum O} 
\usepackage{microtype}

\usepackage{graphicx}
	\graphicspath{{/Users/goby/pictures/teach/163/lecture/}
	{/Users/goby/pictures/teach/common/}} % set of paths to search for images

\usepackage{multicol}
%\usepackage{booktabs}
%\usepackage{textcomp}

\usepackage{tikz}
	\tikzstyle{every picture}+=[remember picture,overlay]
%	\usetikzlibrary{arrows}

\newcommand{\blurred}[3]{
\draw [black!#1] node at (#2,#2) {#3};
\draw [black!#1] node at (0,#2) {#3};
\draw [black!#1] node at (#2,0) {#3};
\draw [black!#1] node at (-#2,#2) {#3};
\draw [black!#1] node at (-#2,0) {#3};
\draw [black!#1] node at (0,-#2) {#3};
\draw [black!#1] node at (#2,-#2) {#3};
\draw [black!#1] node at (-#2,-#2) {#3};
}

\newcommand{\blurry}[1]{
\blurred{5}{0.04}{#1}
\blurred{13}{0.017}{#1}
\blurred{10}{0.012}{#1}
\blurred{15}{0.007}{#1}
\blurred{20}{0.003}{#1}
\blurred{10}{0.0015}{#1}
\draw [black!40] node at (0,0) {#1};
}

\mode<presentation>
{
  \usetheme{Lecture}
  \setbeamercovered{invisible}
  \setbeamertemplate{items}[square]
}

%\usefonttheme[onlymath]{serif}
%\usecolortheme[named=blue7]{structure}


%%% Creative Commons Licenses. Establish, then add to Beamer template.
%\newcommand{\ccbysa}[1][4]{\textsc{cc by-sa #1.0}} % Use version 4.0 as default.
\newcommand{\ccby}[1]{%
	\ifx&#1&
	{\textsc{cc by}}%
\else
	{\textsc{cc by #1.0}}
\fi}


\newcommand{\ccbysa}[1]{%
	\ifx&#1&
	{\textsc{cc by-sa}}%
\else
	{\textsc{cc by-sa #1.0}} 
\fi}

\newcommand{\ccbyncsa}[1]{%
	\ifx&#1&
	{\textsc{cc by-nc-sa}}%
\else
	{\textsc{cc by-nc-sa #1.0}}
\fi}

\newcommand{\ccbync}[1]{%
	\ifx&#1&
	{\textsc{cc by-nc}}%
\else
	{\textsc{cc by-nc #1.0}}
\fi}


\begin{document}

\lecture{student}{student}
\begin{frame}{Our goal for this lecture is to }
	
	\hangpara explain \highlight{genetic drift,} 

	\hangpara and other ways that populations evolve,
	
	\hangpara discuss what are species.
	
\end{frame}
%
%
\begin{frame}[t,plain]{Let's revisit a question.}
	\vspace{2ex}
	
	\onslide<1->\hangpara What is \highlight{evolution?} \onslide<2->\alt<handout>{}{\\Evolution is the genetic change in a population over time.}\\[2\baselineskip]
	

	\onslide<2->\hangpara What does ``genetic change'' mean?

\end{frame}


\begin{frame}{Evolution occurs when there is}
	
	\hangpara natural selection,

	\hangpara gene flow,
	
	\hangpara mutation,
	
	\hangpara non-random mating, and

	\hangpara genetic drift.


	
\end{frame}
%

%{
%\usebackgroundtemplate{\includegraphics[width=\paperwidth]{planktonic_starfish_geneflow} }
%\begin{frame}[b]{}
%\textcolor{white}{\tiny \textcopyright Richard R. Kirby, All Rights Reserved.}
%%Juvenile starfish, \textit{Luidia sarsi}, 
%\end{frame}
%}
%
\begin{frame}[t]{\highlight{Gene flow} is the movement of individuals among populations.}
	\alt<handout>{}{\includegraphics<1>[width=\textwidth]{gene_flow_beetles_start}}
	\alt<handout>{}{\includegraphics<2>[width=\textwidth]{gene_flow_beetles_middle}}
	\includegraphics<3>[width=\textwidth]{gene_flow_beetles_end}

	\onslide<3>When individuals move, their alleles go with them.
	
	\vfilll
	
	\hfill \tiny Tsaneda, Wikimedia, \ccby{3}
\end{frame}
%
%{
%\usebackgroundtemplate{\includegraphics[width=\paperwidth]{mutation} }
%\begin{frame}[b]{}
%\hfill\textcolor{white}{\tiny Illustrator unknown.}
%\end{frame}
%}
%
{
\usebackgroundtemplate{\includegraphics[width=\paperwidth]{mutation_types} }
\begin{frame}[t]{A \highlight{mutation} is any change to the genome of an organism.}

\vfilll

\hfill\rotatebox{90}{\tiny \textcopyright Roberts \& Company}
\end{frame}
}
\begin{frame}[t]{Mutations can create new alleles.}

	\hangpara Most mutations are \highlight{neutral mutations.}\pause
	
	\hangpara Mutations that occur in a gene create new alleles.\pause
	
	\hangpara Some new alleles can be detrimental or beneficial.\pause

	\hangpara More alleles increase \highlight{heterozygosity} in a population.\pause
	
	\hangpara More heterozygosity means more genetic variation for natural selection to work with.
	
\end{frame}
%%
%{
%\usebackgroundtemplate{\includegraphics[width=\paperwidth]{mating_lions} }
%\begin{frame}[b]{}
%\tiny\textcolor{white}{Frank Vassen, Wikimedia, \ccby{2}.}
%\end{frame}
%}
%
{
\usebackgroundtemplate{\includegraphics[width=\paperwidth]{mating_mice} }
\begin{frame}[b]{\highlight{Non-random mating} occurs when individuals choose mates based on phenotype.}
\hfill \tiny \textcopyright Pearson Education, Inc.
\end{frame}
}
%
{
\usebackgroundtemplate{\includegraphics[width=\paperwidth]{charles_ancestry_inbred} }
\begin{frame}[b]{Non-random mating \emph{always} reduces heterozygosity.}
\end{frame}
}
%
%{
%\usebackgroundtemplate{\includegraphics[width=\paperwidth]{snowgeese_drift} }
%\begin{frame}[b]{}
%\end{frame}
%}
\begin{frame}[t]{\highlight{Genetic drift} is change in allele frequencies due to random events.}
	\vspace*{-\baselineskip}
	\begin{center}
		\includegraphics[width=0.7\textwidth]{genetic_drift}
	\end{center}

	\vspace*{-\baselineskip}
	
	\hangpara Genetic drift is strongest in small populations.
	

\end{frame}
{
\usebackgroundtemplate{\includegraphics[width=\paperwidth]{snowgeese_drift_blurred} }
\begin{frame}{Genetic drift \emph{always} causes a \highlight{loss of heterozygosity.}}
	\includegraphics[width=\textwidth]{loss_of_heterozygosity}
\end{frame}
}


{
\usebackgroundtemplate{\includegraphics[width=\paperwidth]{snowgeese_drift_blurred} }
\begin{frame}{Two types of drift cause \highlight{rapid} loss of heterozygosity.}

	\begin{multicols}{2}
		\includegraphics[width=0.5\textwidth]{loss_of_heterozygosity}
	\columnbreak
		\hangpara \highlight{Bottleneck effect} and\\ 
		\highlight{Founder effect.}
	\end{multicols}
	
\end{frame}
}
%
{
\usebackgroundtemplate{\includegraphics[width=\paperwidth]{bottleneck} }
\begin{frame}[b]{\highlight{Bottlenecks} occur when population size decreases rapidly.}

\hfill \tiny \textcopyright Pearson Education, Inc.
\end{frame}
}

\lecture{instructor}{instructor}
{
\usebackgroundtemplate{\includegraphics[width=\paperwidth]{pacific_ocean} }
\begin{frame}[b]{}
\end{frame}
}

{
\usebackgroundtemplate{\includegraphics[width=\paperwidth]{pingelap_atoll} }
\begin{frame}[b]{}
\hfill \tiny \textcolor{white}{Google Earth}
\end{frame}
}

{
\usebackgroundtemplate{\includegraphics[width=\paperwidth]{tropical_island_normal} }
\begin{frame}[b]{}
\tiny Benedict Adam, Flickr, \ccby{2}.
\end{frame}
}



{
\usebackgroundtemplate{\includegraphics[width=\paperwidth]{tropical_island_blurred} }
\begin{frame}[b]{}
\begin{tikzpicture}
\blurry{\tiny\textcolor{gray}{Benedict Adam, Flickr, \ccby{2}.}}
\end{tikzpicture}
\end{frame}
}

\lecture{student}{student}
{
\usebackgroundtemplate{\includegraphics[width=\paperwidth]{pingelap_atoll} }
\begin{frame}[b]{\textcolor{white}{1 in 17 people on Pingelap have achromatopsia.}}
\hfill \tiny \textcolor{white}{Google Earth}
\end{frame}
}

{
\usebackgroundtemplate{\includegraphics[width=\paperwidth]{founder_effect_silver_eye} }
\begin{frame}[b,plain]{}
\tiny \textcopyright Pearson Education \hfill Toby Hudson, Wikimedia, \ccbysa{3}
\end{frame}
}

% Subdividing evolution
\begin{frame}{Evolution can be divided into three related processes.}

	\hangpara \alert<1,5,7>{Microevolution:} evolution within populations.
	\pause

	\hangpara \alert<2,6>{Speciation:} the evolution of new species.
	\pause

	\hangpara \alert<3,8>{Macroevolution:} evolution above the species level.
	\pause

	\centering
	\begin{tikzpicture}
		[remember picture, overlay,
		myLine/.style={color=black,very thick},
		myArrow/.style={color=orange6, thick, ->}]

		\visible<4>{
			\draw [myLine] (-3,-2) -- (0,-2);
			\draw [myLine] (0,-2) -- (0,-0.7); % vertical lines
			\draw [myLine] (0,-2) -- (0,-3.3);
			\draw [myLine] (0,-0.7) -- (3,-0.7); % descendants
			\draw [myLine] (0,-3.3) -- (3,-3.3); 
		}

		% Draw the tree
		\visible<5->{
			\draw [myLine] (-3,-2) -- (0,-2);
		} % ancestor

		\visible<6->{
			\draw [myLine] (0,-2) -- (0,-0.7); % vertical lines
			\draw [myLine] (0,-2) -- (0,-3.3);
		}
		
		\visible<7->{
			\draw [myLine] (0,-0.7) -- (3,-0.7); % descendants
			\draw [myLine] (0,-3.3) -- (3,-3.3); 
		}

		% Ancestor microevolution
		\visible<5>{
			\node [color=orange6] at (-1.5,-1.8) {Microevolution};
		} % ancestor
	
		% Speciation
		\visible<6>{
			\node [color=orange6] (speciation) at (-1.5,-3.1) {Speciation};
			\node [circle, draw=orange6, thick, minimum size=3mm] (speccirc) at (0,-2) {};
			\draw [color=orange6, thick] (speciation) edge (speccirc);
		}
		
		% Descendant microevolution
		\visible<7>{
			\node [color=orange6] at (1.5,-0.5) {Microevolution}; % descendants
			\node [color=orange6] at (1.5,-3.1) {Microevolution};
		}

		% Macroevolution
		\visible<8>{
			\node [color=orange6] (macro) at (1.5,-2) {Macroevolution};
			\draw [myArrow] (macro.north) to (1.5,-0.8);
			\draw [myArrow] (macro.south) to (1.5,-3.2);
		}

	\end{tikzpicture}


%	\vspace{\baselineskip}
%	\hangpara Speciation and macroevolution require microevolution to occur.
\end{frame}

\end{document}
